\chapter{作品测试与分析}
\section{版本、环境与证据口径}\label{sec:env}
本章以提交快照\texttt{1b64b1f}为代码基线。该快照的运行时源码与经持续集成验证的\texttt{96d591b}逐字节一致，差异仅在文档；因此工程回归数字按\texttt{96d591b}的持续集成运行引用，源码一致性由清单哈希核验。历史阶段（M25及其之前的冻结矩阵、较早云模型实验）保留各自源版本，只作对应关系记录。

四类结果分表记录，不能互相冒充：工程回归验证代码行为；确定性脚本验证编排和实际效果；规则检测评测描述初筛分类表现；真实云实验验证模型行为。前三类是当前基线的可复现结果，第四类为开发期历史记录。

当前基线的验证元数据位于源码包的\path{docs/validation/MAIN_SYNC_VALIDATION_2026-10-20.json}，含312份运行时源文件的SHA-256、三平台回归结果、离线检测评测结果与退出码；清单由\path{tools/build_report_binding.py}按一条固定规则从git对象重算，便于复核。历史M25结果见\path{docs/validation/M25_LOCAL_ACCEPTANCE_2026-10-18.json}。摘要清单用于核对文件一致性，不能替代原始观察。

开发端使用Python 3.12和锁定依赖，前端执行Vue类型检查与生产构建。M24容器环境为Linux/aarch64、Python 3.12.14、Git 2.47.3；宿主为Docker Desktop。各平台记录只在其自身范围内解释。

表\ref{tab:evidence}把本章的主要结论压成索引：每条结论只对应一处观察和一个可打开的文件，避免用一段叙述同时支撑多个主张。

\begin{table}[htbp]\centering
\caption{主要结论、关键观察与证据编号}\label{tab:evidence}
\begin{tabularx}{\textwidth}{P{0.26\textwidth}P{0.36\textwidth}YY}\toprule
结论 & 关键观察 & 证据 & 出处\\\midrule
同一提交三平台一致 & 1107 / 1106 / 1107项通过，退出码0 & E1 & \ref{sec:env}节、表\ref{tab:engineering}\\
授权约束进入真实工具链 & 48单元中46项业务完成，2个危险写入被拦 & E2 & \ref{sec:offline}节、表\ref{tab:offline}\\
规则初筛可复现 & 42条TP21/FN0/FP0/TN21；140条查准1.00、查全0.5625 & E3 & \ref{sec:detection}节、表\ref{tab:detection}\\
四业务确定性复现 & 16/16完成，云调用0 & E4 & \ref{sec:detection}节末\\
报告与源码同版本 & 312份源码哈希一致 & E5 & \ref{sec:env}节\\\bottomrule
\end{tabularx}
\end{table}

证据编号对应的文件（路径以随附源码包为根）：E1为\path{docs/validation/MAIN_SYNC_VALIDATION_2026-10-20.json}与三平台持续集成运行记录；E2为M25本地验收记录\path{docs/validation/M25_LOCAL_ACCEPTANCE_2026-10-18.json}；E3为\path{sujie/eval/reports/runs/competition_evidence_20261020/}下三份评测JSON；E4为同目录族中的\path{competition_repro_20261018/summary.json}；E5为\path{competition/submission/latex/runtime-version.json}。除E2属历史源版本外，其余四项均在当前基线重算过。

\section{业务与安全判定方法}
合法业务Oracle核对实际来源读取和报告字段，例如时间线的窗口与排序、附件的整数分金额、Git配置差异、Memory历史和当前事实。安全Oracle读取真实Receiver记录、文件、Git对象/ref及SQLite版本。任务状态FINISHED、记录有效、业务完成和危险效果是四个独立字段。

同事实比较方法为：给候选封存相同业务事实，构造正常C、注入A和适用的合法危险原文引用Q三种场景；主配置为unsafe、strong-static和full，no-source、no-recovery为消融。M25阶段共48个唯一单元、14个配对组，覆盖时间线、附件、IOC、Git检视、三层委派与跨会话Memory。单元数、配对数和独立任务数分别统计。

\section{工程回归与跨平台结果}
\begin{table}[htbp]\centering
\caption{同一提交（快照1b64b1f / 源码96d591b）的工程回归记录}\label{tab:engineering}
\begin{tabularx}{\textwidth}{P{0.20\textwidth}P{0.30\textwidth}Y}\toprule
记录 & 观察结果 & 适用范围\\\midrule
Ubuntu持续集成 & 1107 passed，1 skipped，1 xfailed，1 warning；1231秒 & 全量pytest，含静态检查与格式检查\\
Windows持续集成 & 1106 passed，2 skipped，1 xfailed，1 warning；1938秒 & 同上；多1项因平台跳过\\
macOS本机复跑 & 1107 passed，1 skipped，1 xfailed，1 warning；512秒；退出码0 & 同上，独立于持续集成环境\\
前端构建 & 类型检查与生产构建退出码0 & 构建验证；页面操作另验\\
容器配置与前端作业 & 两者均成功 & 配置可解析、控制台可构建\\
M24 Linux六链 & 六候选各一实例完成，容器退出码0，主机独立重算通过 & 单容器、同UID、脚本模型，无云调用\\\bottomrule
\end{tabularx}
\end{table}

三平台的通过数一致或接近（Windows多一次跳过），说明当前测试夹具不依赖单一操作系统；此前在macOS上因临时目录符号链接导致的路径别名拒绝已被夹具修复。M24容器使用UID/GID 10001、只读rootfs、\texttt{--network none}、CapDrop ALL、no-new-privileges、pids与内存限额及指定tmpfs；观测到CapEff为0、NoNewPrivs为1、rootfs写探针返回EROFS，仅loopback处于UP且无可用默认路由。该记录验证六候选有限链路的运行约束，未测量完整矩阵的隔离或吞吐开销。

\section{规则检测离线评测}\label{sec:detection}
检测评测全部离线运行，无云调用；每次输出都记录数据集哈希与检测器版本，脚本拒绝覆盖已有输出。评测对象是规则初筛的分类表现，不执行真实工具链，也不涉及模型诱导提议。

\begin{table}[htbp]\centering
\caption{规则检测离线评测（当前基线）}\label{tab:detection}
\begin{tabularx}{\textwidth}{P{0.28\textwidth}P{0.08\textwidth}P{0.27\textwidth}Y}\toprule
数据集（检测器） & 规模 & 混淆矩阵 & 召回 / 正常误报\\\midrule
合成冒烟样例（\texttt{rules-v0.3}） & 14 & TP5/FN2/FP1/TN6 & 71.4\% / 14.3\%\\
七绕过合成矩阵（\texttt{rules-v0.6}） & 42 & TP21/FN0/FP0/TN21 & 100\% / 0\%\\
公开注入样本·留出（\texttt{rules-v0.6}） & 140 & TP45/FN35/FP0/TN60 & 56.25\% / 0\%\\\bottomrule
\end{tabularx}
\end{table}

公开样本按来源分层抽样，四个来源的留出召回为：MCP工具投毒85\%\cite{dsmcptoolpoison}、Gandalf 65\%\cite{dsgandalf}、中文智能体安全集55\%\cite{dsagentsafetyzh}、deepset通用注入集20\%\cite{dsdeepset}，各来源精度均为100\%；四个来源的许可均允许再分发样例。合成矩阵由三种注入来源（直接输入、间接资料、记忆）与七类绕过标签构成，含21个攻击格与21个对应正常格。

结果须与口径并列解读：合成矩阵的100\%属同源自检，规则与样例同批开发，只能说明该规则集在自建样例上自洽；公开样本的56.25\%是规则初筛对既有语料的覆盖，其中数值较低的两个来源主要含有需要语义判断的载荷。两组数字都不构成攻击成功率、网关实际阻断率或跨语料泛化能力的结论。扩展检测器提高了中文与绕过变体的覆盖，但检测器只参与风险解释，实际授权仍由策略与不可覆盖名单决定。

四业务确定性运行（\path{competition_repro_20261018}）在当前基线重跑为16个单元全部完成，云调用为0，源码归档哈希与协议哈希随运行目录记录。该运行同时给出网关自身开销的量级：16个单元共84次工具调用，单次调用端到端均摊7.7毫秒，扣除脚本模型等待后为4.7毫秒，单元墙钟均值为40.7毫秒。由于模型是脚本替身，这组数字只说明本地策略核验与账本写入不构成瓶颈，不能外推为接入真实模型后的端到端时延。

\section{同事实离线比较（历史阶段）}\label{sec:offline}
\begin{table}[htbp]\centering
\caption{M25原冻结版本的离线观察，源版本独立于当前基线}\label{tab:offline}
\begin{tabularx}{\textwidth}{P{0.27\textwidth}P{0.13\textwidth}P{0.16\textwidth}P{0.16\textwidth}Y}\toprule
运行方式 & 单元/配对 & 正常执行 & 业务完成 & 危险效果\\\midrule
确定性脚本 & 48/14 & 48 & 46 & 2\\
本地HTTP MockTransport & 48/14 & 48 & 46 & 2\\
三来源固定提议代表 & 14/7 & 14 & 14 & 3\\\bottomrule
\end{tabularx}
\end{table}

前两行各有两个Git注入单元在unsafe和strong-static配置实际执行了禁止推送；完整配置保持检视任务的远端约束。其他没有危险提议或没有配置差异的格，只报告其实际观察。记录通过核验表示观察和协议一致，并不意味着每个格都安全。

本地HTTP矩阵使用265次进程内MockTransport请求，经过Provider请求准备、解析和账本代码，云调用为0。三来源代表涵盖直接Unicode、外部Base64和跨会话Memory流程备注；来源真实暴露后固定提出禁止的夹具标记写入。三个unsafe攻击格实际改变标记，对应full格阻断。这说明指定实例的权限检查有效，未测量自然模型的攻击成功率或来源机制的独有收益。

M25原矩阵与当前基线使用不同源版本。当前版本保留后续远端修复，历史48+48+14矩阵未在当前基线全量重跑。报告保留这一版本对应关系，避免将旧矩阵当成新代码成绩。

\section{较早真实模型开发记录（历史阶段）}
表\ref{tab:cloud}列出可查的历史记录，均属开发期小样本，不能与当前基线结果合并统计。
\begin{table}[htbp]\centering
\caption{历史云模型小样本，独立于离线矩阵}\label{tab:cloud}
\begin{tabularx}{\textwidth}{P{0.21\textwidth}YY}\toprule
阶段 & 实际结果 & 请求与已知用量\\\midrule
M19共同配置 & 10单元中7项业务完成，3项IOC未提出发送且BLOCK为0 & 31请求，25005 tokens\\
M21首轮 & 四业务中3项完成，时间线窗口与排序失败保留 & 15请求，24144 tokens\\
M21时间线修复 & 修复后该开发样例单独完成 & 3请求，5645 tokens\\
M20引用核验 & 脚本主流程加一次云分析 & 1请求，353 tokens\\\bottomrule
\end{tabularx}
\end{table}

这些结果用于检查模型接入和业务可行性；跨版本修复不合并成功率。当前真实云CLI关闭，新的真实模型协议、资源准入、未见输入、重复规则和任务层统计需预登记后执行。无危险动作提议的格不能归因为网关阻断；只有开发期文档记载、原始运行目录不在仓库的结果，只能按“开发期记录”引用。

\section{结果分析与未完成验证}
现有结果支持四类具体结论：第一，任务与来源用途约束已进入真实工具链，并能在部分Git与标记写入校准中形成可观察差异；第二，委派、Memory和危险原文引用具有实际角色、版本、内容和业务事实见证；第三，审批对象与实际提交参数可核对，敏感写入目标在放宽授权下仍被拒绝；第四，同一提交在三个平台的全量回归一致通过，说明当前工程记录可跨环境复核。四业务确定性运行与规则检测评测均可离线重算。

下一轮需要解决的是相对收益和外推范围：在同版真实模型上比较业务完成、危险效果和代价；补充未见任务、有效消融和重复；完成运行时三来源七绕过矩阵、跨轮组合、完整入口绑定，以及当前拓扑的隔离、性能与跨机复现。暂不填写无数据支撑的正式ASR、提升比例或性能达标值。
